%! title: 梯度、策略梯度与 PPO
%! date: 2026-10-25
%! slug: rl-ppo
%! series: 强化学习
%! number: 1
%! tags: 强化学习, PPO, 策略梯度, GAE
%! description: 从导数与梯度出发，推导策略梯度、优势函数、Actor-Critic、重要性采样与 GAE，最后落到 TRPO 与 PPO 的 clip 目标。
\documentclass[UTF8,fontset=fandol]{ctexart}

% 版式
\usepackage{geometry}
\geometry{
  a4paper,
  margin=2.5cm,
  headheight=14pt,
}

% 主题色（克制、印刷友好）
\usepackage{xcolor}
\definecolor{themeaccent}{HTML}{1a5f7a}
\definecolor{themecolorthm}{HTML}{1a5f7a}
\definecolor{themecolordef}{HTML}{2d6a4f}
\definecolor{themecolorremark}{HTML}{8b6914}

% 定理块底纹浓度（xcolor：数字越大底色越深；旧版约 !4–!6，此处加深便于对比）
\newcommand{\thmshadecolor}{themecolorthm!18}
\newcommand{\defshadecolor}{themecolordef!14}
\newcommand{\remarkshadecolor}{themecolorremark!12}

% 数学
\usepackage{amsmath,amssymb}
\usepackage{mathtools}
\usepackage{amsthm}
\usepackage[anycase]{thmtools}
\usepackage{bm}

% 行距（可读性）
\usepackage{setspace}
\onehalfspacing

% 页眉页脚
\usepackage{fancyhdr}
\pagestyle{fancy}
\fancyhf{}
\fancyhead[L]{\small\color{themeaccent!70!black}\nouppercase{\leftmark}}
\fancyhead[R]{\small\thepage}
\renewcommand{\headrulewidth}{0.4pt}
\renewcommand{\headrule}{\color{themeaccent!50!black}\hrule height 0.4pt}

% 章节标题着色 + 装饰线（ctex）
\ctexset{
  section/format     = \Large\bfseries\color{themeaccent},
  section/aftername  = \hspace{0.5em},
  section/aftertitle = \par\vspace{-4pt}%
    {\color{themeaccent!40}\rule{\linewidth}{0.5pt}}\vspace{4pt},
  subsection/format  = \large\bfseries\color{themeaccent!88!black},
}

% 常用配对符号（mathtools）
\DeclarePairedDelimiter{\paren}{(}{)}
\DeclarePairedDelimiter{\sqbrk}{[}{]} % 勿用 \brack，与内核冲突
\DeclarePairedDelimiter{\curlbrace}{\{}{\}} % 勿用 \brace，与内核冲突
\DeclarePairedDelimiter{\abs}{\lvert}{\rvert}
\DeclarePairedDelimiter{\norm}{\lVert}{\rVert}

% 定理样式：彩色标题（底纹在 \declaretheorem 中用 shaded 指定）
\declaretheoremstyle[
  spaceabove=6pt,
  spacebelow=6pt,
  headfont=\bfseries\color{themecolorthm},
  notefont=\bfseries\color{themecolorthm!80!black},
  bodyfont=\itshape,
  headpunct={.},
  postheadspace=0.5em,
]{styletheorem}

\declaretheoremstyle[
  spaceabove=6pt,
  spacebelow=6pt,
  headfont=\bfseries\color{themecolordef},
  notefont=\bfseries\color{themecolordef!80!black},
  bodyfont=\normalfont,
  headpunct={.},
  postheadspace=0.5em,
]{styledefinition}

\declaretheoremstyle[
  spaceabove=6pt,
  spacebelow=6pt,
  headfont=\bfseries\color{themecolorremark},
  notefont=\bfseries\color{themecolorremark!80!black},
  bodyfont=\normalfont,
  headpunct={.},
  postheadspace=0.5em,
]{styleremark}

\declaretheorem[
  style=styletheorem,
  name=定理,
  numberwithin=section,
  shaded={bgcolor=\thmshadecolor, margin=0pt},
]{theorem}

\declaretheorem[
  style=styletheorem,
  name=引理,
  sibling=theorem,
  shaded={bgcolor=\thmshadecolor, margin=0pt},
]{lemma}

\declaretheorem[
  style=styletheorem,
  name=命题,
  sibling=theorem,
  shaded={bgcolor=\thmshadecolor, margin=0pt},
]{proposition}

\declaretheorem[
  style=styletheorem,
  name=推论,
  sibling=theorem,
  shaded={bgcolor=\thmshadecolor, margin=0pt},
]{corollary}

\declaretheorem[
  style=styledefinition,
  name=定义,
  sibling=theorem,
  shaded={bgcolor=\defshadecolor, margin=0pt},
]{definition}

\declaretheorem[
  style=styleremark,
  name=注,
  sibling=theorem,
  shaded={bgcolor=\remarkshadecolor, margin=0pt},
]{remark}

\renewcommand{\proofname}{\textcolor{themecolorthm}{证明}}
\renewcommand{\qedsymbol}{\textcolor{themecolorthm}{$\blacksquare$}}

% 彩色标题页（替代 \maketitle）
\makeatletter
\newcommand{\maketitlecolored}{%
  \begin{center}\leavevmode
    {\LARGE\bfseries\color{themeaccent}\@title\par}
    \vspace{0.8em}
    {\large\@author\par}
    {\normalsize\color{black!55}\@date\par}
  \end{center}
  \vspace{1.8em}%
}
\makeatother

% 表格（专业三线表）
\usepackage{booktabs}

% 图表标题样式
\usepackage{caption}
\captionsetup{labelfont={bf,color=themeaccent}, labelsep=period, font=small}

% 链接（置于 preamble 末尾）
\usepackage{hyperref}
\hypersetup{
  colorlinks=true,
  linkcolor=themeaccent!90!black,
  citecolor=themecolordef!75!black,
  urlcolor=themeaccent,
}

% 智能交叉引用（须在 hyperref 之后加载）
\usepackage{cleveref}
\crefname{theorem}{定理}{定理}
\crefname{lemma}{引理}{引理}
\crefname{proposition}{命题}{命题}
\crefname{corollary}{推论}{推论}
\crefname{definition}{定义}{定义}
\crefname{remark}{注}{注}
\crefformat{equation}{式~(#2#1#3)}

\title{强化学习梯度与ppo笔记}
\author{THT}
\date{2026.10.25}

\begin{document}
\maketitlecolored

\section{梯度与导数}

\subsection{导数}

函数 $f$ 在点 $x$ 处的导数定义为
\[
f'(x)=\lim_{h\to 0}\frac{f(x+h)-f(x)}{h}
\]
其中 $\dfrac{f(x+h)-f(x)}{h}$ 是曲线上两点 $(x,f(x))$ 与 $(x+h,f(x+h))$ 连线（割线）的斜率。当 $h\to 0$ 时，割线趋于切线，所以 $f'(x)$ 就是曲线在 $x$ 处切线的斜率，也就是函数在该点的变化率。

当 $h$ 很小时，
\[
f(x+h)\approx f(x)+f'(x)\,h
\]
即在 $x$ 附近，曲线可以用切线近似：输入变化 $h$，输出大约变化 $f'(x)\,h$。

只有一个自变量时，只能沿 $x$ 轴向左或向右移动：$f'(x)>0$ 时向右函数增大，$f'(x)<0$ 时向左函数增大，$|f'(x)|$ 表示变化的快慢。自变量变多后，各个方向的变化率组成一个向量，即梯度。

\subsection{多元函数：偏导数组成梯度}

对二元函数 $f(x,y)$，$\dfrac{\partial f}{\partial x}$ 是只沿 $x$ 方向移动时的变化率，$\dfrac{\partial f}{\partial y}$ 是只沿 $y$ 方向移动时的变化率。同时移动一小步 $(\Delta x,\Delta y)$ 时，
\[
\Delta f\approx\frac{\partial f}{\partial x}\Delta x+\frac{\partial f}{\partial y}\Delta y=\nabla f\cdot(\Delta x,\Delta y)
\]
其中
\[
\nabla f=\left(\frac{\partial f}{\partial x},\frac{\partial f}{\partial y}\right)
\]
就是梯度。几何上，函数在该点附近可以用切平面近似，梯度描述了这个切平面的倾斜方向和陡峭程度。

\subsection{方向导数：梯度方向上升最快}

沿单位向量 $\vec{u}$（$\|\vec{u}\|=1$）移动时，函数的变化率（方向导数）为
\[
D_{\vec{u}}f=\nabla f\cdot\vec{u}=\|\nabla f\|\,\|\vec{u}\|\cos\theta=\|\nabla f\|\cos\theta
\]
其中 $\theta$ 是 $\vec{u}$ 与 $\nabla f$ 的夹角。因此：
\begin{itemize}
  \item $\theta=0$（沿梯度方向）：$\cos\theta=1$，变化率取最大值 $\|\nabla f\|$，函数上升最快；
  \item $\theta=\pi/2$（垂直于梯度，即沿等高线）：变化率为 $0$；
  \item $\theta=\pi$（沿负梯度方向）：函数下降最快，这就是梯度下降沿 $-\nabla f$ 更新的原因。
\end{itemize}
这里的“最快”是局部的：在当前点沿各个方向走同样长的一小步，沿梯度方向函数值增加最多。走远之后梯度会改变，所以梯度上升（下降）每次只走一小步，并在新位置重新计算梯度。

\section{强化学习方法}

\subsection{Policy-based}

在给定状态下，策略给出各个动作的概率。我们要优化策略，使它更倾向于选择好的动作，从而获得更高的回报。

\subsubsection{符号定义}

\begin{itemize}
  \item $\tau=(s_0,a_0,s_1,a_1,\dots,s_T)$：一条轨迹。
  \item $\pi_\theta(\cdot\mid s_t)$：随机策略，即状态 $s_t$ 下动作的概率分布；$\pi_\theta(a_t\mid s_t)$ 是选择动作 $a_t$ 的概率，$\theta$ 为策略参数。
  \item $P(s_{t+1}\mid s_t,a_t)$：状态转移概率，即在 $s_t$ 执行 $a_t$ 后转移到 $s_{t+1}$ 的概率。它由环境决定，与 $\theta$ 无关。
  \item $r_t=r(s_t,a_t,s_{t+1})$：在 $s_t$ 执行 $a_t$ 并转移到 $s_{t+1}$ 时获得的即时奖励。
  \item $R(\tau)$：轨迹的累计折扣奖励（回报），$\gamma\in(0,1]$ 为折扣因子：
  \[
  R(\tau)=\sum_{t=0}^{T-1}\gamma^t r_t
  \]
  \item $P(\tau\mid\pi_\theta)$：在策略 $\pi_\theta$ 下产生轨迹 $\tau$ 的概率，$\rho_0(s_0)$ 为初始状态是 $s_0$ 的概率：
  \[
  P(\tau\mid\pi_\theta)=\rho_0(s_0)\prod_{t=0}^{T-1}\pi_\theta(a_t\mid s_t)\,P(s_{t+1}\mid s_t,a_t)
  \]
\end{itemize}

\subsubsection{优化目标}

策略的好坏用期望回报衡量：
\[
J(\pi_\theta)=\mathbb{E}_{\tau\sim\pi_\theta}\big[R(\tau)\big]=\sum_{\tau}P(\tau\mid\pi_\theta)\,R(\tau)
\]
$J(\pi_\theta)$ 是按策略 $\pi_\theta$ 采样轨迹时回报的期望，即所有轨迹回报按概率加权的平均。优化目标是找到使它最大的参数：
\[
\theta^*=\arg\max_{\theta}J(\pi_\theta)
\]

\subsubsection{策略梯度}

用梯度上升求解，需要计算
\[
\nabla_\theta J(\pi_\theta)=\sum_{\tau}R(\tau)\,\nabla_\theta P(\tau\mid\pi_\theta)
\]
直接计算有两个困难：$P(\tau\mid\pi_\theta)$ 是很多项的连乘，并且含有未知的环境转移概率；对所有轨迹求和也无法实现。

利用对数求导技巧
\[
\nabla_\theta\log f(\theta)=\frac{\nabla_\theta f(\theta)}{f(\theta)}
\quad\Longrightarrow\quad
\nabla_\theta f(\theta)=f(\theta)\,\nabla_\theta\log f(\theta)
\]
可得
\begin{align*}
\nabla_\theta J(\pi_\theta)
&=\sum_{\tau}R(\tau)\,P(\tau\mid\pi_\theta)\,\nabla_\theta\log P(\tau\mid\pi_\theta)\\
&=\mathbb{E}_{\tau\sim\pi_\theta}\big[R(\tau)\,\nabla_\theta\log P(\tau\mid\pi_\theta)\big]
\end{align*}
第二步由期望的定义得到。这样，对所有轨迹求和变成了期望，可以用采样来估计。

再对轨迹概率取对数，连乘变为连加：
\begin{align*}
\log P(\tau\mid\pi_\theta)
&=\log\rho_0(s_0)+\sum_{t=0}^{T-1}\log P(s_{t+1}\mid s_t,a_t)\\
&\quad+\sum_{t=0}^{T-1}\log\pi_\theta(a_t\mid s_t)
\end{align*}
前两项与 $\theta$ 无关，梯度为 $0$，所以
\[
\nabla_\theta\log P(\tau\mid\pi_\theta)=\sum_{t=0}^{T-1}\nabla_\theta\log\pi_\theta(a_t\mid s_t)
\]
代入上式，得到策略梯度：
\[
\nabla_\theta J(\pi_\theta)=\mathbb{E}_{\tau\sim\pi_\theta}\left[R(\tau)\sum_{t=0}^{T-1}\nabla_\theta\log\pi_\theta(a_t\mid s_t)\right]
\]
环境转移概率在求导时被消去，因此不需要知道环境模型，也能计算策略梯度。

\subsubsection{小结}

基于优化目标 $J(\pi_\theta)$，策略梯度为
\[
\nabla_\theta J(\pi_\theta)=\mathbb{E}_{\tau\sim\pi_\theta}\left[\sum_{t=0}^{T-1}R(\tau)\,\nabla_\theta\log\pi_\theta(a_t\mid s_t)\right]
\]
实际中采样 $N$ 条轨迹，用样本均值估计这个期望（无偏，但方差很大）：
\[
\nabla_\theta J(\pi_\theta)\approx\frac{1}{N}\sum_{n=0}^{N-1}\sum_{t=0}^{T_n-1}R(\tau_n)\,\nabla_\theta\log\pi_\theta\big(a_t^{n}\mid s_t^{n}\big)
\]
其中 $s_t^{n},a_t^{n}$ 是第 $n$ 条轨迹 $\tau_n$ 在时刻 $t$ 的状态和动作，$T_n$ 是它的长度。然后沿梯度上升更新参数（$\alpha$ 为学习率）：
\[
\theta\leftarrow\theta+\alpha\,\nabla_\theta J(\pi_\theta)
\]

\subsection{价值函数}

上式用整条轨迹的回报 $R(\tau)$ 给每个动作打分，方差较大。更一般地，用 $\Psi_t$ 表示衡量动作 $a_t$ 好坏的权重，它可以是回报，也可以是价值函数。梯度可以写成
\begin{align*}
\nabla_\theta J(\pi_\theta)
&=\mathbb{E}_{\tau\sim\pi_\theta}\left[\sum_{t=0}^{T-1}\Psi_t\,\nabla_\theta\log\pi_\theta(a_t\mid s_t)\right]\\
&\approx\frac{1}{N}\sum_{n=0}^{N-1}\sum_{t=0}^{T_n-1}\Psi_t^{n}\,\nabla_\theta\log\pi_\theta\big(a_t^{n}\mid s_t^{n}\big)
\end{align*}
$\Psi_t$ 常见的形式如下（$Q_\pi$、$A_\pi$、$V_\pi$ 的定义见下文）：
\begin{center}
\renewcommand{\arraystretch}{1.4}
\begin{tabular}{cll}
\hline
 & $\Psi_t$ & 含义\\
\hline
1 & $\sum_{t'=0}^{T-1}\gamma^{t'}r_{t'}$ & 整条轨迹的回报 $R(\tau)$\\
2 & $\sum_{t'=t}^{T-1}\gamma^{t'-t}r_{t'}$ & 动作 $a_t$ 之后的回报（reward-to-go）\\
3 & $\sum_{t'=t}^{T-1}\gamma^{t'-t}r_{t'}-b(s_t)$ & 上一项减去基线 $b(s_t)$\\
4 & $Q_\pi(s_t,a_t)$ & 动作价值函数\\
5 & $A_\pi(s_t,a_t)$ & 优势函数\\
6 & $r_t+\gamma V_\pi(s_{t+1})-V_\pi(s_t)$ & TD 误差（TD residual）\\
\hline
\end{tabular}
\end{center}

\subsubsection{状态价值函数}

记从时刻 $t$ 开始的回报为
\[
G_t=\sum_{k=0}^{T-t-1}\gamma^k r_{t+k}=r_t+\gamma G_{t+1}
\]
状态价值函数是从状态 $s_t$ 出发、按策略 $\pi$ 行动的期望回报：
\[
V_\pi(s_t)=\mathbb{E}_\pi\left[G_t\mid s_t\right]
\]
展开得
\begin{align*}
V_\pi(s_t)
&=\mathbb{E}_\pi\left[r_t+\gamma G_{t+1}\mid s_t\right]\\
&=\mathbb{E}_\pi\left[r_t\mid s_t\right]+\gamma\,\mathbb{E}_\pi\left[V_\pi(s_{t+1})\mid s_t\right]\\
&=\sum_{a_t\in\mathcal{A}}\pi(a_t\mid s_t)\sum_{s_{t+1}\in\mathcal{S}}P(s_{t+1}\mid s_t,a_t)\,r(s_t,a_t,s_{t+1})\\
&\quad+\gamma\sum_{a_t\in\mathcal{A}}\pi(a_t\mid s_t)\sum_{s_{t+1}\in\mathcal{S}}P(s_{t+1}\mid s_t,a_t)\,V_\pi(s_{t+1})\\
&=\sum_{a_t\in\mathcal{A}}\pi(a_t\mid s_t)\sum_{s_{t+1}\in\mathcal{S}}P(s_{t+1}\mid s_t,a_t)\big[r_t+\gamma V_\pi(s_{t+1})\big]\\
&=\mathbb{E}_{a_t\sim\pi(\cdot\mid s_t)}\Big[\mathbb{E}_{s_{t+1}\sim P(\cdot\mid s_t,a_t)}\big[r_t+\gamma V_\pi(s_{t+1})\big]\Big]
\end{align*}
第二步利用全期望公式和马尔可夫性，把 $\mathbb{E}_\pi\left[G_{t+1}\mid s_t\right]$ 写成了 $\mathbb{E}_\pi\left[V_\pi(s_{t+1})\mid s_t\right]$。

\subsubsection{动作价值函数}

动作价值函数是在状态 $s_t$ 执行动作 $a_t$、之后按策略 $\pi$ 行动的期望回报：
\[
Q_\pi(s_t,a_t)=\mathbb{E}_\pi\left[G_t\mid s_t,a_t\right]
\]
展开得
\begin{align*}
Q_\pi(s_t,a_t)
&=\mathbb{E}_\pi\left[r_t+\gamma G_{t+1}\mid s_t,a_t\right]\\
&=\mathbb{E}_\pi\left[r_t\mid s_t,a_t\right]+\gamma\,\mathbb{E}_\pi\left[V_\pi(s_{t+1})\mid s_t,a_t\right]\\
&=\sum_{s_{t+1}\in\mathcal{S}}P(s_{t+1}\mid s_t,a_t)\big[r(s_t,a_t,s_{t+1})+\gamma V_\pi(s_{t+1})\big]\\
&=\mathbb{E}_{s_{t+1}\sim P(\cdot\mid s_t,a_t)}\big[r_t+\gamma V_\pi(s_{t+1})\big]
\end{align*}
与 $V_\pi$ 相比，动作 $a_t$ 已经给定，所以只需对下一状态 $s_{t+1}$ 求期望。

\subsubsection{状态价值与动作价值的关系、优势函数}

比较两个展开式可知，$V_\pi$ 就是 $Q_\pi$ 对动作的期望：
\[
V_\pi(s_t)=\mathbb{E}_{a_t\sim\pi(\cdot\mid s_t)}\big[Q_\pi(s_t,a_t)\big]=\sum_{a_t\in\mathcal{A}}\pi(a_t\mid s_t)\,Q_\pi(s_t,a_t)
\]
优势函数衡量动作 $a_t$ 比当前策略的平均水平好多少：
\[
A_\pi(s_t,a_t)=Q_\pi(s_t,a_t)-V_\pi(s_t)
\]
由于 $V_\pi(s_t)$ 与 $s_{t+1}$ 无关，可以放进期望中：
\begin{align*}
A_\pi(s_t,a_t)
&=\mathbb{E}_{s_{t+1}\sim P(\cdot\mid s_t,a_t)}\big[r_t+\gamma V_\pi(s_{t+1})\big]-V_\pi(s_t)\\
&=\mathbb{E}_{s_{t+1}\sim P(\cdot\mid s_t,a_t)}\big[r_t+\gamma V_\pi(s_{t+1})-V_\pi(s_t)\big]\\
&=\mathbb{E}_{s_{t+1}\sim P(\cdot\mid s_t,a_t)}\big[\delta_t\big]
\end{align*}
其中 $\delta_t=r_t+\gamma V_\pi(s_{t+1})-V_\pi(s_t)$ 称为 TD 误差（TD error）。当价值函数准确时，$\delta_t$ 是优势 $A_\pi(s_t,a_t)$ 的无偏估计。

\subsection{Actor-Critic}

Actor 是策略网络 $\pi_\theta$，负责选择动作；Critic 是价值网络 $V_\phi$，负责评估状态价值，并为 Actor 提供优势的估计。

\subsubsection{Actor 的更新}

用 Critic 给出的 TD 误差作为优势的估计：
\[
\hat{A}_t=r_t+\gamma V_\phi(s_{t+1})-V_\phi(s_t)
\]
代入策略梯度（即取 $\Psi_t=\hat{A}_t$）：
\begin{align*}
\nabla_\theta J(\pi_\theta)
&\approx\mathbb{E}_{\tau\sim\pi_\theta}\left[\sum_{t=0}^{T-1}\hat{A}_t\,\nabla_\theta\log\pi_\theta(a_t\mid s_t)\right]\\
&\approx\frac{1}{N}\sum_{n=0}^{N-1}\sum_{t=0}^{T_n-1}\hat{A}_t^{n}\,\nabla_\theta\log\pi_\theta\big(a_t^{n}\mid s_t^{n}\big)
\end{align*}
实现时通常写成按单步平均的目标函数，并把 $\hat{A}_t$ 当作常数（不对它求导）：
\[
L^{\text{actor}}(\theta)=\mathbb{E}_t\big[\hat{A}_t\log\pi_\theta(a_t\mid s_t)\big]
\]
它对 $\theta$ 的梯度就是上面的策略梯度估计，对它做梯度上升（或对 $-L^{\text{actor}}$ 做梯度下降）即可。注意它只是用来求梯度的替代目标，数值上并不等于 $J(\pi_\theta)$。

\subsubsection{Critic 的更新}

Critic 的目标是让预测值 $V_\phi(s_t)$ 逼近目标值（TD target）。

\paragraph{优化目标}
使用均方误差（MSE）作为损失，求使它最小的 $\phi$：
\[
L(\phi)=\mathbb{E}_t\Big[\big(\underbrace{r_t+\gamma V_\phi(s_{t+1})}_{\text{TD target}}-\underbrace{V_\phi(s_t)}_{\text{预测值}}\big)^2\Big]
\]
括号内就是 TD 误差 $\delta_t=r_t+\gamma V_\phi(s_{t+1})-V_\phi(s_t)$。注意 TD target 本身也是估计值，并不是真实价值。

\paragraph{梯度}
如果对 TD target 也求导，得到完整梯度：
\[
\nabla_\phi L=\mathbb{E}_t\Big[2\delta_t\big(\gamma\nabla_\phi V_\phi(s_{t+1})-\nabla_\phi V_\phi(s_t)\big)\Big]
\]
实际中把 TD target 当作常数（detach），只对预测值求导，称为半梯度（semi-gradient）：
\[
\nabla_\phi L\approx\mathbb{E}_t\big[-2\delta_t\,\nabla_\phi V_\phi(s_t)\big]
\]
这是 TD 学习的标准做法：目标值不随本次更新移动，训练更稳定。

\paragraph{直观理解}
梯度下降的更新为（$\beta$ 为学习率）
\[
\phi\leftarrow\phi-\beta\,\nabla_\phi L=\phi+2\beta\,\delta_t\,\nabla_\phi V_\phi(s_t)
\]
\begin{itemize}
  \item $\delta_t>0$（预测偏低）：$V_\phi(s_t)$ 被\textbf{调高}；
  \item $\delta_t<0$（预测偏高）：$V_\phi(s_t)$ 被\textbf{调低}。
\end{itemize}

\paragraph{PyTorch 伪代码}\mbox{}
\begin{verbatim}
values = critic_net(states)
with torch.no_grad():  # TD target 不参与求导
    next_values = critic_net(next_states)
    # dones=1 表示终止状态，没有后续价值
    td_target = rewards + gamma * next_values * (1 - dones)

critic_loss = nn.MSELoss()(values, td_target)
critic_optimizer.zero_grad()  # 清空上一步的梯度
critic_loss.backward()
critic_optimizer.step()
\end{verbatim}

\subsubsection{Actor-Critic 的迭代过程}

\paragraph{策略改进（Actor）}
设 Critic 能准确评估当前策略 $\pi$（即 $V_\phi=V_\pi$），按优势更新策略参数：
\[
\theta\leftarrow\theta+\alpha\,A_\pi(s_t,a_t)\,\nabla_\theta\log\pi_\theta(a_t\mid s_t)
\]
$A_\pi(s_t,a_t)>0$ 时提高 $\pi_\theta(a_t\mid s_t)$，$A_\pi(s_t,a_t)<0$ 时降低，由此得到新策略 $\pi'$。

\paragraph{何时达到最优}
记 $s_t$ 下最好的动作为 $a_t^*=\arg\max_a Q_\pi(s_t,a)$。利用 $\sum_a\pi(a\mid s_t)=1$：
\begin{align*}
A_\pi(s_t,a_t^*)
&=Q_\pi(s_t,a_t^*)-\sum_{a\in\mathcal{A}}\pi(a\mid s_t)\,Q_\pi(s_t,a)\\
&=\sum_{a\in\mathcal{A}}\pi(a\mid s_t)\big[Q_\pi(s_t,a_t^*)-Q_\pi(s_t,a)\big]\\
&=\sum_{a\neq a_t^*}\underbrace{\pi(a\mid s_t)}_{\text{选其他动作的概率}}\,
\underbrace{\big[Q_\pi(s_t,a_t^*)-Q_\pi(s_t,a)\big]}_{\text{动作间的差距}\,\geq\,0}
\end{align*}
当 $\pi(a_t^*\mid s_t)\to 1$ 时，其他动作的概率趋于 $0$，于是
\[
V_\pi(s_t)\to Q_\pi(s_t,a_t^*),\qquad A_\pi(s_t,a_t^*)\to 0
\]
若每个状态都如此，策略已无改进空间，即达到最优策略 $\pi^*$。

注意：优势趋于 $0$ 并不表示动作分不出好坏。方括号中的差距始终存在，趋于 $0$ 的是选其他动作的概率；较差动作的优势 $A_\pi(s_t,a)\to Q_\pi(s_t,a)-Q_\pi(s_t,a_t^*)<0$，只是它们几乎不再被采样。

\paragraph{价值评估（Critic）}
策略由 $\pi$ 更新为 $\pi'$ 后，关系式 $V(s_t)=\sum_a\pi(a\mid s_t)\,Q(s_t,a)$ 右侧的动作分布随之改变，旧的 $V_\pi$ 不再适用，需要重新拟合：
\[
V_{\pi'}(s_t)=\sum_{a\in\mathcal{A}}\pi'(a\mid s_t)\,Q_{\pi'}(s_t,a)
\]
实际中通过最小化前面的 Critic 损失 $L(\phi)$ 完成这一步。当 $V_\phi=V_\pi$ 时，$\mathbb{E}\left[\delta_t\mid s_t,a_t\right]=A_\pi(s_t,a_t)$，所以 critic loss 趋于 $0$ 与优势趋于 $0$ 相对应。

\paragraph{小结}
Critic 评估当前策略，Actor 依据优势改进策略，二者交替进行。当优势与 critic loss 都趋于 $0$ 时：
\begin{itemize}
  \item Actor：每个状态下 $\pi_\theta(a_t^*\mid s_t)\to 1$，即 $\pi_\theta\to\pi^*$；
  \item Critic：准确评估当前策略，$V_\phi\to V_\pi$；随着策略逼近最优，$V_\phi\to V_{\pi^*}$。
\end{itemize}

\subsection{重要性采样与GAE}
对于Actor的梯度优化目标而言，
\begin{align*}
\nabla_\theta J(\pi_\theta)
&\approx\mathbb{E}{\tau\sim\pi\theta}\left[\sum_{t=0}^{T-1}\hat{A}t\,\nabla\theta\log\pi_\theta(a_t\mid s_t)\right]\\
&\approx\frac{1}{N}\sum_{n=0}^{N-1}\sum_{t=0}^{T_n-1}\hat{A}t^{n}\,\nabla\theta\log\pi_\theta\big(a_t^{n}\mid s_t^{n}\big)
\end{align*}
在实际运行中，我们每次针对策略\(\pi_\theta\)采样后，得到相应的数据\(r_t,V_\phi(s_{t+1},\nabla log\pi_\theta(a_t|s_t)\)等，再根据这些数据，带入公式并更新参数，参数更新后，在旧策略\(\pi_{\theta_{old}}\)下采样的数据就不能使用了，需要重新采样数据，再由重新采样的数据更新参数（on-policy），而每一次重新采样会带来较大的资源消耗，由于每次采样的随机性很大，此时，采样带来的方差往往很大；更重要的是，在实际训练过程中，用critic网络拟合出来\(V_\pi\)并不一定是准确的价值函数，所以，我们使用的TD error去估计优势其实是有偏差的。

为此，我们希望同一条旧数据能用很多次，研究者引入：

重要性采样（important Sampling）：让旧数据也能被反复利用几次
GAE：平衡优势函数的方差和偏差
经验回放：把以前的数据存起来，像抽盲盒一样混合着抽出来用（off-policy）

重要性采样：
为了降低采样成本，我们希望采样一批的数据能用来更新神经网络多次，
假设某次更新完毕后，我们得到策略\(\pi_{old}\)
我们用\(\pi_{old}\)和环境交互，得到一批回合数据
我们把这批回合数据重复使用k次：我们先把这批数据喂给\(\pi_{old}\)得到\(\pi_{\theta_1}\)，再把这批数据喂给\(\pi_{\theta_1}\)得到\(\pi_{\theta_2}\),以此类推，当重复k次后，我们得到了\(\pi_\theta\)

\begin{itemize}
\item 假设有两个分布 \(p(x), q(x)\)
\item 最开始我想从 \(p(x)\) 中进行多次采样，然后求函数 \(f(x)\) 的期望。例如我想从 \(\pi_\theta\) 中进行采样，然后求累积奖励的期望，这个期望我们表示成 \(E_{x \sim p(x)}[f(x)]\)
\item 但是现在，因为某些原因，我们无法从 \(p(x)\) 中直接采样，只能从另一个分布 \(q(x)\) 中进行采样了，那么此时我们要怎么表示 \(E_{x \sim p(x)}[f(x)]\)？
\item 为了解决这个问题，我们做如下变换：
\end{itemize}



\begin{align*}
E_{x \sim p(x)} [f(x)] &= \int p(x) f(x) dx \\
&= \int \frac{p(x)}{q(x)} q(x) f(x) dx \\
&= E_{x \sim q(x)} \left[ \frac{p(x)}{q(x)} f(x) \right]
\end{align*}


但在实际操作中，我们会遇到p(x)和q(x)分布差异较大的问题，这个时候，我们需要足够多的采样来抵消这种差异对期望的最终影响。

\textbf{重要性采样前，策略的梯度是：}

\[
\nabla J(\pi_\theta) = E_t [A_\phi(s_t, a_t) \nabla \log \pi_\theta(a_t|s_t)]
\]

\textbf{重要性采样后，策略的梯度是：}

\[
\nabla J(\pi_\theta) = \underset{\tau \sim \pi_{\theta_{old}}}{E_t} \left[ \frac{\pi_\theta(a_t|s_t)}{\pi_{old}(a_t|s_t)} A_\phi(s_t, a_t) \nabla \log \pi_\theta(a_t|s_t) \right]
\]

\textbf{我们根据重要性采样构造了这个新的策略梯度，那么对应的新的 actor 优化目标就可以从这个策略梯度中反推出来：}

\[
\arg \max_{\pi_\theta} J(\pi_\theta) = \underset{\tau \sim \pi_{\theta_{old}}}{E_t} \left[ \frac{\pi_\theta(a_t|s_t)}{\pi_{\theta_{old}}(a_t|s_t)} A_\phi(s_t, a_t) \right]
\]

GAE：平衡优势函数的方差和偏差
在前面我们假设了\(V_\pi\)能正确评估策略\(\pi\)的情况下，我们利用\(r_t + \gamma V_\pi(s_{t+1}) - V_\pi(s_t)\)作为优势函数\(A_\pi\)的无偏估计，但在训练过程中，\(V_\pi\) 往往无法完全准确评估策略\(\pi\)的价值。

\textbf{为了解决因为 \(V_\pi\) 估计不准而引发的“高偏差”问题，直观上我们可以尽量少信任 \(V_\pi\) 的策略，即对于 \(r_t + \gamma V_\pi(s_{t+1}) - V_\pi(s_t)\)，我们可以把 \(V_\pi(s_{t+1})\) 做递归地展开，得到：}

\[
-V_\pi(s_t) + \sum_{l=0}^{\infty} \gamma^l r_{t+l}
\]

其中，\(r_t, r_{t+1}, r_{t+2}, \dots\) 都是我们某次采样得到的即时奖励数据。如果 \(V_\pi\) 不准，那么我就信任我的实际采样结果，这样至少不会让我对优势函数的估计出现偏差。

但采取这种做法又会引发一个新问题：我们知道 \(r_t, r_{t+1}, r_{t+2}, \dots\) 它们都是随机变量，相比之前只用 \(r_t\)，现在的做法带来的随机性更大了（相当于在每一个 timestep 都引入了随机性，随机性逐步累加）。也就是如果之前的方差是 \(Var(r_t)\)，那么现在的方差则变成更大。

当 \(V_\pi\) 估计不准时，如果我能在这 \(r_t + \gamma V_\pi(s_{t+1}) - V_\pi(s_t)\) 和 \(-V_\pi(s_t) + \sum_{l=0}^{\infty} \gamma^l r_{t+l}\) 间取得一种平衡，即我既不想完全信任 \(V_\pi\)，又不想完全信任我全部的采样的结果 \(r_t, r_{t+1}, r_{t+2}, \dots\)，\textbf{此时，GAE 通过一个超参数 \(\lambda\) 控制了我们想要的方差-偏差间的平衡。}

\[
\Psi_t = \sum_{l=0}^{\infty} (\gamma\lambda)^l \delta_{t+l}
\]

\begin{itemize}
\item \(\delta_t = r_t + \gamma V_\pi(s_{t+1}) - V_\pi(s_t)\)
\item \(\gamma\)：超参，折扣因子
\item \(\lambda\)：\textbf{超参，即用于平衡方差-偏差的因子。}
\begin{itemize}
\item 当\(\lambda\)接近0时，\(\Psi_t\)退化成 \(r_t + \gamma V_\pi(s_{t+1}) - V_\pi(s_t)\)，也即上图所绘制的高偏差情况
\item 当\(\lambda\)接近1时，\(\Psi_t\)变成 \(-V_\pi(s_t) + \sum_{l=0}^{\infty} \gamma^l r_{t+l}\)，也即上图所绘制的高方差-低偏差情况
\item \textbf{综上，\(\lambda\)越小，方差越小，偏差越大。\(\lambda\)越大，方差越大，偏差越小。}
\end{itemize}
\end{itemize}

\textbf{在接下来的表达中，我们记这种引入了GAE方法的单步优势为} \(A_\phi^{GAE}(s_t, a_t)\)

在引入重要性采样和GAE后，我们的优化目标变成\[
arg \max_{\pi_\theta} J(\pi_\theta) = \underset{\tau \sim \pi_{\theta_{old}}}{E_t} \left[ \frac{\pi_\theta(a_t|s_t)}{\pi_{\theta_{old}}(a_t|s_t)} A_\phi^{GAE}(s_t, a_t) \right]
\]

\subsection{TRPO}
在Actor-critic框架中，对于citic，我们的优化目标为\[\min_\phi L(\phi) = E_t \left[ \left( V_\phi(s_t) - V_{target} \right)^2 \right]\]

对于Actor，我们的优化目标为，\[
arg \max_{\pi_\theta} J(\pi_\theta) = \underset{\tau \sim \pi_{\theta_{old}}}{E_t} \left[ \frac{\pi_\theta(a_t|s_t)}{\pi_{\theta_{old}}(a_t|s_t)} A_\phi^{GAE}(s_t, a_t) \right]
\]

但如果\(\pi_\theta\)和\(\pi_{old}\)这两个分布差异太大，且我们采样的轨迹数量没有足够大时，\(J(\pi_\theta\)的估计是不准确的。

TRPO的解决办法是，把 \(\pi_\theta\) 和 \(\pi_{old}\) 的分布相似性作为 \(J(\pi_\theta)\) 的 constraint，\textbf{这就是 TRPO 的做法，即我们有：}

\begin{align*}
arg \max_{\pi_\theta} J(\pi_\theta) &= \underset{\tau \sim \pi_{\theta_{old}}}{E_t} \left[ \frac{\pi_\theta(a_t|s_t)}{\pi_{old}(a_t|s_t)} A_\phi^{GAE}(s_t, a_t) \right] \\[1.5ex]
subject\ to\ E_t &[KL(\pi_{\theta_{old}}(\cdot|s_t), \pi_\theta(\cdot|s_t))] \leq \delta
\end{align*}

由于这种限制不直接加在 \(J(\pi_\theta)\) 中，因此使得整体优化过程变得较为复杂，

\subsection{PPO}
我们重新来端详引入重要性采样和GAE以后的强化学习优化目标：

\[
arg \max_{\pi_\theta} J(\pi_\theta) = \underset{\tau \sim \pi_{\theta_{old}}}{E_t} \left[ \frac{\pi_\theta(a_t|s_t)}{\pi_{\theta_{old}}(a_t|s_t)} A_\phi^{GAE}(s_t, a_t) \right]
\]

\begin{itemize}
\item \textbf{当 \(A_\phi^{GAE}(s_t, a_t) > 0\) 时，说明我们当前执行的动作 \(a_t\) 相较于别的动作更好，因此我们要提升 \(\pi_\theta(a_t|s_t)\)。}但是考虑到在采样不足的情况下，\(\pi_\theta\) 和 \(\pi_{\theta_{old}}\) 的分布差异不能太大，因此 \(\frac{\pi_\theta(a_t|s_t)}{\pi_{\theta_{old}}(a_t|s_t)}\) 是有上限的，即不能一味轻信 \(A_\phi^{GAE}(s_t, a_t)\) 而持续不断提升 \(\pi_\theta(a_t|s_t)\)，你要保证 \(\pi_\theta\) 在 \(\pi_{\theta_{old}}\) 的信任阈内。\textbf{我们假设当这个比值 \(>= 1 + \epsilon\) 的时候，就对这个比值做 clip，固定在 \(1 + \epsilon\)。}
\item \textbf{当 \(A_\phi^{GAE}(s_t, a_t) < 0\) 时，说明我们当前执行的动作 \(a_t\) 相较于别的动作更差，因此我们要降低 \(\pi_\theta(a_t|s_t)\)。}考虑到 \(\pi_\theta\) 和 \(\pi_{old}\) 的分布差异不能太大，因此 \(\frac{\pi_\theta(a_t|s_t)}{\pi_{\theta_{old}}(a_t|s_t)}\) 是有下限的，即不能一味轻信 \(A_\phi^{GAE}(s_t, a_t)\) 而持续不断降低 \(\pi_\theta(a_t|s_t)\)，\textbf{我们假设当这个比值 \(<= 1 - \epsilon\) 的时候，就对这个比值做 clip，固定在 \(1 - \epsilon\)。}
\item \textbf{\(\epsilon\) 就是 PPO 当中的一个超参。}
\end{itemize}

基于这种 clip 思想，我们得出了 ppo-clip 下强化学习的优化目标，为了表达简便，我们这里令 \(r_t(\theta) = \frac{\pi_\theta(a_t|s_t)}{\pi_{\theta_{old}}(a_t|s_t)}\)：

\[
J^{CLIP}(\pi_\theta) = \underset{\tau \sim \pi_{\theta_{old}}}{E_t} \left[ min \left[ r_t(\theta)A_\phi^{GAE}(s_t, a_t), \ clip(r_t(\theta), 1-\epsilon, 1+\epsilon)A_\phi^{GAE}(s_t, a_t) \right] \right]
\]


在PPO中，critic网络的loss设计为：
\begin{align*}
V_t^{CLIP} &= clip(V_t^{new}, V_t^{old} - \epsilon, V_t^{old} + \epsilon) \\
R_t &= A_t^{GAE} + V_t^{old}
\end{align*}

\[
arg \min_{V_\phi} L(V_\phi) = E_t \left[ max \left[ (V_t^{new} - R_t)^2, (V_t^{CLIP} - R_t)^2 \right] \right]
\]

我们来详细解释这个公式的含义：

\begin{itemize}
\item 首先，我们先根据 8.3 (2) 中对 GAE 的定义，递归地展开 \(A_t^{GAE}\)，得到：
\[
A_t^{GAE} = \delta_t + \gamma\lambda A_{t+1}^{GAE} = (r_t + \gamma * V_{t+1}^{old} - V_t^{old}) + \gamma\lambda A_{t+1}^{GAE}
\]

\item 基于这个表达式，本质上 \(R_t = A_t^{GAE} + V_t^{old} = (r_t + \gamma * V_{t+1}^{old}) + \gamma\lambda A_{t+1}^{GAE}\)

\item 在 ppo epoch 开始前，我们有初始 actor 和初始 critic。其中初始 critic 就是上面所说的 \(V_t^{old}\)

\item 我们用初始 actor 和初始 critic 产生经验值，然后这波经验值将被用在接下来 ppo epochs 轮的迭代中，也就是在这些迭代中，\textbf{上式里的 \(r_t, V_t^{old}, V_{t+1}^{old}, A_t^{GAE}, A_{t+1}^{GAE}, R_t\) 都来自这些初始的经验值，它们在接下来的这个 ppo epochs 轮迭代中是不变的。}

\item \textbf{本质上，critic 的优化目标是以下差值的平方：} \\
\(V_t^{new} - R_t = [V_t^{new} - (r_t + \gamma * V_{t+1}^{old})] - \gamma\lambda A_{t+1}^{GAE}\)。这里我们先不考虑 clip，因为 clip 只是一种 trick。
\end{itemize}

\begin{itemize}
\item \textbf{在理解了这些基本点的基础上，我们再来考虑引入 clip 的作用，它和 actor 的 ppo-clip 一样，}本质上是想让 critic 在信任阈内更新：
\begin{itemize}
\item 当 \(V_t^{new} < R_t\) 时，我们确实应该鼓励 \(V_t^{new}\) 向 \(R_t\) 提升，但是如果它本身的值已经超出 \(V_t^{old}\) 的信任阈 \(V_t^{old} + \epsilon\) 之外，此时我们就应该对它做 clip。也就是说，如果此时我们应用 max，最终取的是 \((V_t^{CLIP} - R_t)^2\)，注意，\(V_t^{CLIP}\) 最终是由 \(V_t^{old}\) 定义的，而这一项来自我们之前收集的经验，是个独立的常数，它和等待更新的这个 critic 无关，所以相当于我们本次不对 critic 做梯度更新。
\item 当 \(V_t^{new} > R_t\) 时，也是同理类推，这里不赘述。
\end{itemize}
\end{itemize}

\section{参考}
1、https://www.zhihu.com/people/lemonround/posts

2、https://www.xiaohongshu.com/explore/6addc9fe00000000250057e0?app_platform=android\&ignoreEngage=true\&app_version=9.49.1\&share_from_user_hidden=true\&xsec_source=app_share\&type=normal\&xsec_token=CBTDga-aRbhPGpX2PUUzTnA-kAeTzBeo5DCf0fU57KXsU=\&author_share=1\&xhsshare=WeixinSession\&shareRedId=N0o1Mjw1SkI2NzUyOTgwNjc1OTg1Okg-\&apptime=1792925623\&share_id=d7541f253f7a4b3592c4d44159bf7dd0\&share_channel=wechat\&track_c

\end{document}
